\section{Discussion}
\label{sec:discussion}
The mean life-expectancy cost to the individual donor ($\approx 909$~days, or 2.5~years) must be weighed against the benefit conferred on the recipient and, through the voucher mechanism, on designated family members who may one day develop ESRD. Contemporary studies estimate kidney transplant candidates experience a gain of five life years from transplantation,~\cite{stedman2025estimating} and the ESRD-conditional model in this work demonstrates that priority voucher program is worth 109~days to a family member who develops ESRD at age~60. Viewed in this frame, donation represents a transfer of life expectancy from donor to recipient on favorable terms.

This aggregate figure conceals some heterogeneity. Age at donation is the demographic factor that matters most: the life-expectancy cost is roughly 91\% larger at age~25 relative to age~55, because a younger donor lives long enough after the roughly decade-long delay in donor mortality risk for a much larger share of their remaining life to fall under the elevated hazard ratio. Race and sex differ by about 7\% and 3\%, respectively, at the base-case donation age. This analysis supports age-, race-, and sex-specific life-expectancy disclosure: a 25-year-old Black male and a 55-year-old White female donor are, in effect, being asked to accept quantitatively different bargains under the same generic counseling framework.

The model also reframes what priority-access policy is actually protecting. KAS250 and the NKR voucher mechanism are designed around a rare event, an eventual ESRD diagnosis in a prior donor or designated family member, and this analysis shows that their unconditional, population-level life-expectancy value is correspondingly small because so few beneficiaries ever draw on it. This is not evidence that such policies are unimportant: conditional on the rare event they are designed for, priority access can be worth hundreds of days at the individual level, and the dilution reflects low incidence and effective screening rather than low value per event. The appropriate way to communicate these mechanisms may therefore be closer to catastrophic insurance, whose expected value is small precisely because it is rarely invoked, than to an intervention expected to move average outcomes.

The results are dominated by a single structural uncertainty, the time-course of the donor all-cause mortality hazard ratio (HR), rather than by sampling error in any individually estimated parameter. Short-to-medium-follow-up cohorts and the meta-analysis pooling them report no excess mortality, while the one cohort followed long enough to plausibly detect a late-emerging effect reports a hazard ratio of 1.30 that its own authors attribute to risk emerging only after roughly a decade. This analysis encodes that pattern as a piecewise-linear ramp (flat at 1.0 through year~10, rising to 1.30 by year~15, matching the Mjøen et~al.\ cohort's median follow-up), which is a minimal, directly literature-motivated functional form but an assumption nonetheless. The true magnitude and time-course could differ. If the true HR is closer to 1.0, race, sex, and age play far greater roles.

% The model reports only undiscounted life-years and does not incorporate quality-of-life weighting; because dialysis is associated with substantially reduced health utility, a quality-adjusted life year (QALY)-adjusted analysis would likely widen, not narrow, the donor's disadvantage during any dialysis-dependent interval.

% Several parameters were necessarily borrowed across source populations under an explicit sourcing rule (e.g., within-donor hazard ratios from Massie et al.\ applied to the Muzaale matched-control baseline), and no single cohort traces the complete pathway from living donation through ESRD onset to transplantation, so the calibration checks reported here constitute face validity rather than external validation against a purpose-built dataset.

I therefore join Mjøen's call\cite{mjoeen_reply} for for Garg,\cite{garg2012} Segev,\cite{segev_perioperative_2010} and their colleagues to ``publish updated results from their study cohorts,'' which could now span 22~years of follow-up, becoming the longest large-scale studies of donor all-cause mortality ever conducted. Future work should prioritize prospective donor cohorts with follow-up long enough, ideally beyond 15--20 years, to directly estimate whether and when donor all-cause mortality risk begins to rise, since this single question now determines whether donation carries a life-expectancy cost at all.